\documentclass[11pt]{article}

\usepackage[review]{acl}

\usepackage{times}
\usepackage{latexsym}
\usepackage[T1]{fontenc}
\usepackage[utf8]{inputenc}
\usepackage{microtype}
\IfFileExists{inconsolata.sty}{\usepackage{inconsolata}}{}    % optional
\usepackage{graphicx}
\usepackage{booktabs}
\usepackage{amsmath}
\usepackage{amssymb}
\usepackage{xcolor}
\usepackage{hyperref}

% Macros
\newcommand{\tool}[1]{\texttt{#1}}
\newcommand{\genetrace}{\textsc{GeneTrace}}
\newcommand{\evoopd}{\textsc{evo-OPD}}
\newcommand{\geneexam}{\textsc{GENE-bench}}
\newcommand{\ideaevolving}{\textsc{IdeaEvolving}}
\newcommand{\todo}[1]{\textcolor{red}{[TODO: #1]}}
\newcommand{\note}[1]{\textcolor{blue}{[note: #1]}}

\title{\genetrace{} and \evoopd{}: A Training-Grade Genome-Centric Corpus and Lineage-Aware Distillation for Scientific Idea Models}

\author{Anonymous Authors \\
  Affiliation \\
  \texttt{email@domain} \\}

\begin{document}
\maketitle

\begin{abstract}
We present \genetrace, the first training-grade corpus instantiating the
\emph{genome-centric} paper-representation paradigm, and \evoopd, a
lineage-aware on-policy distillation algorithm designed to exploit it.
Building on our prior work that defined the paradigm at proof-of-concept
scale (\ideaevolving / \geneexam), \genetrace{} scales the six-field
genome decomposition and five-mode evolutionary-dynamics taxonomy from a
1{,}029-item evaluation set to a 5K--50K corpus annotated with evidence
quotes, automated verifier scores, and teacher reasoning traces. \evoopd{}
extends vanilla on-policy distillation with two new signals --- a
verifier-anchored reward decoupled by per-token role, and a self-supervised
lineage-consistency signal --- both of which generalise to open-ended
generation. On \geneexam{} main challenge, fine-tuning Qwen3-8B with
\genetrace-derived data and \evoopd{} reaches \todo{X}\% accuracy, a
\todo{Y}-point improvement over an SFT-only baseline at matched compute.
We release the corpus, training code, contamination guard, and trained
model weights under permissive licenses.
\end{abstract}

\section{Introduction}
\label{sec:intro}

Recent progress in scientific-idea generation by large language models
(\todo{cite ResearchAgent, AutoSurvey, Sakana AI Scientist}) has exposed a
representation gap: current systems reason over papers as either
\emph{document-centric} blobs (abstracts, full text) or as
\emph{method-entity-centric} graphs that link technique names across the
literature (\todo{cite Intern-Atlas}). Neither layer captures \emph{why} one
paper followed another --- whether it fixed a predecessor's limitation,
hybridised two lineages, or speciated into a new niche. This is precisely the
question a model must answer to propose the \emph{next} paper, rather than
merely retrieve a past one.

Our prior work, \ideaevolving~\citep{ideaevolving2025}, argued that a
\emph{genome-centric} layer can fill this gap. Each paper is decomposed into
six fields --- \emph{mechanism}, \emph{niche}, \emph{observation},
\emph{limitation}, \emph{delta}, \emph{claim} --- and each lineage edge
between two papers is tagged with one of five \emph{evolutionary dynamics}:
\textit{Mutation}, \textit{Adaptive Radiation}, \textit{Hybridization},
\textit{Speciation}, or \textit{Niche Competition}. \ideaevolving{}
demonstrated this paradigm by hand-curating the 1{,}029-item \geneexam{}
evaluation set and showing that frontier models struggle on it
(\todo{cite GPT-5.5 number}).

This paper takes the next step: turning the paradigm from an
\textbf{evaluation construct} into a \textbf{training substrate}. We
contribute:

\begin{enumerate}
  \item \textbf{\genetrace{}}: the first release-grade, training-scale
    instantiation of the genome paradigm. Compared to \ideaevolving's eval set,
    \genetrace{} adds per-field evidence quotes, automated verifier scores on
    every annotation, and teacher reasoning traces, and scales from $\sim$2K
    hand-curated papers to 5K--50K papers annotated by a strong teacher
    (GPT-5.5) under explicit contamination guards. We ship the denylist and
    safe-pool construction code with the release, enabling provable disjoint
    splits between training and \geneexam{} evaluation.
  \item \textbf{\evoopd{}}: a lineage-aware on-policy distillation algorithm
    that uses three signals coupled at the token level: (i) field-gated
    reverse-KL to a strong teacher, (ii) a verifier-anchored reward decoupled
    by per-token role, and (iii) a self-supervised lineage-consistency
    signal derived from \genetrace's lineage chains. The algorithm generalises
    to open-ended generation by replacing the closed-form verifier with a
    continuous one, without altering the optimisation procedure.
  \item \textbf{Empirical demonstration} on Qwen3-8B: \evoopd{} reaches
    \todo{X}\% on \geneexam{} main challenge, \todo{Y} points above an SFT
    baseline and \todo{Z} points above a vanilla-OPD baseline at matched
    compute, with corresponding lifts on the open-ended \texttt{GENE-Arena}
    benchmark.
\end{enumerate}

\paragraph{Why a new training algorithm?} Vanilla on-policy distillation
(Tinker; Qwen3 strong-to-weak; \todo{cite}) treats every token as an
equally-weighted reverse-KL target against a single teacher's distribution.
Verifier-based RL (DAPO, Dr.~GRPO; \todo{cite}) assumes one scalar reward
per trajectory. The genome-centric setting violates both assumptions:
content tokens (the values of the six genome fields) carry the task
information, while boilerplate (JSON brackets, field names) carries none;
and our verifier provides multiple substructure scores rather than a single
exact-match flag. \evoopd{} is the minimal algorithmic modification that
exploits both: per-token role gating, plus a verifier head that decouples
substructure correctness from teacher mimicry.

\paragraph{Why a new corpus?} \ideaevolving's eval-only release does not
suffice for training: the 1{,}029 instances are too few for supervised
fine-tuning, lack evidence-grounding for verifier-based RL, and overlap with
the evaluation set. Concurrent work Intern-Atlas (\todo{cite arXiv
2604.28158}) provides 1M papers at method-entity granularity but lacks the
genome decomposition and dynamics taxonomy required to train a model on the
\emph{why} of paper evolution. \genetrace{} fills the gap between these two:
training-scale (5K--50K), structured (six fields + five dynamics), grounded
(per-field evidence quotes), and reward-ready (per-annotation verifier
scores).

\section{Related Work}
\label{sec:related}

\paragraph{Paper-as-graph corpora.} Several efforts have organised the
scientific literature as a graph. Citation-only graphs (Semantic Scholar
Open Research Corpus \todo{cite}, OpenAlex \todo{cite}, ArXiv) provide
breadth but no semantic edge typing. Concurrent with our work,
\emph{Intern-Atlas} (\todo{cite arXiv 2604.28158}) builds a methodological
evolution graph on 1M papers, with method-level entities and lineage edges
grounded in source text. \genetrace{} is complementary: where Intern-Atlas
covers breadth at method-entity granularity, \genetrace{} covers depth at
full genome-card granularity. Method-entity edges record \emph{that} a
technique evolved; genome-card edges with dynamics labels record
\emph{why}.

\paragraph{The genome paradigm.} Our prior work \ideaevolving{}
\citep{ideaevolving2025} introduced the six-field genome representation and
five-mode evolutionary-dynamics taxonomy. The associated \geneexam{}
benchmark hand-curates 1{,}029 evaluation instances across 42 task types
(T1: paper-level genome extraction; T2: ordering and grouping; T3:
dynamics classification; T4: cross-paper consistency). \ideaevolving{}
demonstrated that the paradigm is human-curatable and that frontier models
score below 25\% on the benchmark. This paper takes the next step: scaling
the paradigm to a training-grade corpus and showing that a small open
model can be trained on it to a non-trivial fraction of frontier
performance.

\paragraph{On-policy distillation.} On-policy distillation (OPD) trains a
student on its own rollouts using token-level reverse-KL to a stronger
teacher. The Tinker recipe \citep{tinker_opd_2025} and the Qwen3 technical
report \citep{qwen3_2025} both demonstrate that OPD recovers Pass@K
behaviour of the teacher while preserving student efficiency. These
recipes use a single signal (teacher distribution) and ignore task
structure. \evoopd{} extends OPD with two additional signals --- a
verifier reward and a lineage-consistency self-signal --- coupled at the
per-token level via role-based gating.

\paragraph{Reinforcement learning from verifiers.} Verifier-based RL
methods such as DAPO \citep{dapo2025} and Dr.~GRPO \citep{drgrpo2025}
replace human-feedback reward models with deterministic verifiers (math
correctness, code unit-tests). These methods provide one scalar reward per
trajectory and are agnostic to teacher distributions or task structure.
\evoopd{} reuses the verifier-reward primitive but applies it at the
per-token-role level, alongside teacher distillation, and adds a
self-supervised lineage signal.

\paragraph{Scientific-idea generation.} Concurrent systems including
ResearchAgent \todo{cite}, AutoSurvey \todo{cite}, and AI Scientist
\todo{cite} use prompt-based pipelines over frontier models to generate
research proposals or surveys. These systems rely on document-centric
context and do not train a specialised model. Our work is complementary:
\genetrace{} could serve as the structured-context layer under such a
pipeline, and a model trained with \evoopd{} can replace the frontier
backbone at a fraction of the compute cost.

\paragraph{Open-ended generation evaluation.} The PES rubric introduced in
\ideaevolving{} (\textit{Heredity}, \textit{Variation}, \textit{Selection}
\citep{ideaevolving2025}) provides a structured judge for open-ended
proposal writing. We adopt PES as the open-ended verifier head inside
\evoopd, integrating it with the schema and evidence checks from the
closed-form verifier.

\section{\genetrace}
\label{sec:genetrace}

\genetrace{} is a four-level annotation corpus over a contamination-guarded
paper pool. Section~\ref{sec:genetrace:schema} defines the schema,
Section~\ref{sec:genetrace:construction} describes the construction
pipeline, and Section~\ref{sec:genetrace:contamination} details the
contamination guard that distinguishes our release from prior paper-graph
corpora.

\subsection{Schema}
\label{sec:genetrace:schema}

\paragraph{Level 1 --- GenomeCard.} For each paper $p$, a JSON record
containing the six gene fields from \ideaevolving{} (mechanism, niche,
observation, limitation, delta, claim), each accompanied by zero or more
\emph{evidence quotes} from the source text. Evidence quotes are
character-offset-anchored and verified by exact substring match against
\texttt{title + abstract + intro}. The full schema is given in
Appendix~\ref{app:schema:card}.

\paragraph{Level 2 --- DynamicsEdge.} For each ordered pair $(p, q)$ in a
lineage chain, a record giving the dynamics label
$d \in \{\text{Mutation}, \text{Adaptive Radiation},
\text{Hybridization}, \text{Speciation}, \text{Niche Competition}\}$, the
\emph{driver} field (which of $p$'s six fields motivated the transition),
and per-field \emph{gene fates}
$\in \{\text{INHERITED}, \text{MUTATED}, \text{LOST}, \text{NOVEL},
\text{HYBRIDIZED}\}$. Edges include a teacher chain-of-thought trace
explaining the dynamics assignment.

\paragraph{Level 3 --- LineageChain.} For each chain of length $\geq 3$, a
record giving the member papers, the per-step dynamics sequence, and a
teacher narrative summary. Chains are constructed by DFS over the
DynamicsEdge graph with cycle detection.

\paragraph{Level 4 --- VerifierBundle.} Per-annotation automated scores
(schema validity, evidence groundedness, dynamics consistency,
fate--dynamics agreement) bundled in a separate file for downstream
consumers who want to filter by quality or use the scores as RL rewards.

\subsection{Construction pipeline}
\label{sec:genetrace:construction}

\paragraph{Paper pool.} We start from $\sim$3,846 pre-2017 papers extracted
from the \ideaevolving{} \texttt{paper\_db}, filtered to those with
abstracts of $\geq$200 words (yielding $\sim$857 in the round-1 SFT pool;
expanded to 5K for the v0.1 release). Pre-2017 papers and the
\ideaevolving{} denylist (Section~\ref{sec:genetrace:contamination}) jointly
guarantee disjointness from \geneexam{} evaluation papers.

\paragraph{Teacher annotation.} For each paper, we prompt GPT-5.5 (Azure
keyless endpoint, frozen snapshot \texttt{2024-12-01-preview}) with the
source text and a structured \emph{gene\_card\_extract} prompt. The teacher
returns a JSON object with the six gene fields, each accompanied by 1--3
verbatim evidence quotes from the source. We accept the annotation iff a
deterministic verifier (Appendix~\ref{app:verifier}) returns
\texttt{schema\_valid} = 1.0 and \texttt{evidence\_grounded} $\geq 0.8$.
For round-1, 100\% of teacher outputs passed the schema check; 856 of 857
passed the evidence check.

\paragraph{Lineage edges.} For each pair $(p, q)$ where $q$ cites $p$ and
both are in the safe pool, we prompt the teacher with both cards and the
candidate dynamics enum. Dynamics validity (label $\in$ enum) and
fate--dynamics consistency (e.g., \textit{Hybridization} requires $\geq 1$
\texttt{HYBRIDIZED} fate) are enforced by the verifier.

\paragraph{Quality filtering.} We discard annotations failing any verifier
check. For the v0.2 release, we additionally apply \emph{consistency
filtering} \citep{west2023symbolic}: each lineage edge is annotated twice
(once by GPT-5.5, once by Qwen3-14B-Thinking), and only edges where both
teachers agree on the dynamics label are retained.

\subsection{Contamination guard}
\label{sec:genetrace:contamination}

A paper-graph corpus intended for training models that will be evaluated on
benchmarks derived from the same literature must explicitly demonstrate
non-overlap with the evaluation set. We adopt three layers:

\begin{enumerate}
  \item \textbf{Denylist.} An explicit blocklist of 8,359 paper IDs
    extracted from \ideaevolving's \texttt{paper\_db}; 90.4\% have a resolvable
    Semantic Scholar ID. We additionally include a 1-hop OpenAlex
    citation expansion (in progress; v0.2 release).
  \item \textbf{Temporal cut.} The training pool is restricted to
    pre-2017 papers. \geneexam{} draws predominantly from 2018+ literature,
    providing a temporal disjointness guarantee independent of the
    denylist.
  \item \textbf{Min-K\%++ leakage diagnostic.} We measure
    \citep{shi2024minkpp}-style memorisation probability that the trained
    model has seen each \geneexam{} paper. We report this number in the
    dataset card; target Min-20\%++ $<$ 0.05 (consistent with no leakage).
\end{enumerate}

The denylist, safe-pool construction code, and Min-K\%++ measurement
script ship with the release, enabling any downstream consumer to verify
the contamination guard end-to-end.

\subsection{Dataset statistics}
\label{sec:genetrace:stats}

\begin{table}[t]
  \centering
  \small
  \begin{tabular}{lrrrr}
    \toprule
    Release & Cards & Edges & Chains & Status \\
    \midrule
    v0.1 (paper) & \todo{5K} & \todo{2K} & \todo{200} & in progress \\
    v0.2 (1mo)   & 20K & 8K & 1K & planned \\
    v1.0         & 50K+ & 25K+ & 5K+ & long-term \\
    \bottomrule
  \end{tabular}
  \caption{\genetrace{} release tiers. v0.1 ships with the paper.}
  \label{tab:release_tiers}
\end{table}

Table~\ref{tab:release_tiers} summarises the tiered release. The v0.1
corpus shipped with this paper contains \todo{5K} GenomeCards
($\sim$\todo{4.3K} with full evidence quotes), \todo{2K} DynamicsEdges
spanning all five modes, and \todo{200} LineageChains of length 3--5.
All annotations include teacher reasoning traces (mean length
\todo{N} tokens, capped at 1024).

\section{\evoopd: Lineage-Aware On-Policy Distillation}
\label{sec:method}

\evoopd{} extends on-policy distillation with two additional signals
required for the genome-centric setting: a \emph{verifier-anchored
reward} that exploits the substructure of \genetrace's verifier, and a
\emph{lineage-consistency self-signal} that uses lineage chains to provide
supervision in the absence of gold answers. Both signals are coupled with
the standard teacher reverse-KL via per-token role gating.

\subsection{Per-token role gating}
\label{sec:method:phi}

For each token $t$ in a student rollout $y$, a deterministic
parser~\citep{evo_opd_parser_impl} assigns one of six role tags
$\phi(t) \in \{$boilerplate, content\_field, evidence\_span,
dynamics\_label, gold\_answer, unknown$\}$ based on the rollout's JSON
structure and the GeneTrace schema. Each role has a fixed weight $\alpha(\phi(t))$
(Table~\ref{tab:phi_weights}): boilerplate (JSON brackets, field names)
receives low weight; content fields and dynamics labels receive high
weight. Token-level gating concentrates the optimisation signal on the
tokens that actually carry task semantics.

\begin{table}[t]
  \centering
  \small
  \begin{tabular}{lc}
    \toprule
    Role $\phi$ & Weight $\alpha(\phi)$ \\
    \midrule
    boilerplate     & 0.3 \\
    content\_field  & 1.0 \\
    evidence\_span  & 1.5 \\
    dynamics\_label & 2.0 \\
    gold\_answer    & 0.0 \\
    unknown         & 0.5 \\
    \bottomrule
  \end{tabular}
  \caption{Per-role weights $\alpha(\phi)$. \texttt{gold\_answer} tokens
  are zero-weighted because they are handled separately by the verifier
  head (Sec.~\ref{sec:method:verifier}).}
  \label{tab:phi_weights}
\end{table}

\subsection{Per-token reward}
\label{sec:method:reward}

For each token $t$, the per-token reward combines three signals:

\begin{align}
  r_t = &- \alpha(\phi(t)) \cdot \mathrm{KL}[\pi_\theta \,\|\, \pi_T](y_t) \nonumber \\
        &+ \alpha(\phi(t)) \cdot \lambda_v \cdot v_{\mathrm{adv}}(y, t) \nonumber \\
        &+ \alpha(\phi(t)) \cdot \lambda_c \cdot c_{\mathrm{adv}}(y, p, t)
  \label{eq:rt}
\end{align}

where $\pi_\theta$ is the student, $\pi_T$ is the teacher,
$v_{\mathrm{adv}}$ is the per-token verifier advantage
(Sec.~\ref{sec:method:verifier}), and $c_{\mathrm{adv}}$ is the
per-token lineage-consistency advantage (Sec.~\ref{sec:method:lineage}).
$\lambda_v$ and $\lambda_c$ are scalar loss weights.

The policy is updated by standard PG on
$\sum_t r_t \cdot \nabla \log \pi_\theta(y_t | y_{<t}, x)$ with no critic
and no reference-KL term (these are subsumed by the teacher-KL
contribution).

\subsection{Verifier-anchored reward}
\label{sec:method:verifier}

The verifier $v: y \mapsto [0, 1]$ from \genetrace{}
(Sec.~\ref{sec:genetrace}) is composed of four subscores:
$v_{\mathrm{schema}}$, $v_{\mathrm{evidence}}$, $v_{\mathrm{dynamics}}$,
$v_{\mathrm{match}}$. For closed-form tasks (e.g., \geneexam{}),
$v_{\mathrm{match}}$ is a 0/1 exact match against gold. For open-ended
tasks, $v_{\mathrm{match}}$ is replaced by a continuous PES rubric judge
\citep{ideaevolving2025} that scores Heredity, Variation, and Selection
on a 0--10 scale and normalises to $[0, 1]$.

The per-token advantage distributes the trajectory-level score across
tokens proportional to their role weight:
\begin{equation}
  v_{\mathrm{adv}}(y, t) = (v(y) - \bar{v}) \cdot \mathbb{1}[\phi(t) \in \mathcal{R}_v]
  \label{eq:vadv}
\end{equation}
where $\mathcal{R}_v = \{$content\_field, evidence\_span, dynamics\_label$\}$
and $\bar{v}$ is a running baseline.

\subsection{Lineage-consistency self-signal}
\label{sec:method:lineage}

The lineage-consistency score $c(y, p) \in [0, 1]$ measures whether the
student's output $y$ on paper $p$ is consistent with $p$'s lineage
context. For tasks where $y$ proposes a successor paper, $c$ is computed
by re-prompting the student to predict the \emph{predecessor} of the
proposed successor and comparing to the ground-truth predecessor card
from \genetrace; high $c$ means the proposal is lineage-faithful. For
tasks where $y$ classifies a relationship between two papers, $c$ is
computed by checking the prediction on a held-out third paper in the
same chain.

Critically, $c$ requires no gold answer for $y$ itself --- only the
lineage structure of the \genetrace{} corpus. This makes it applicable
to open-ended tasks where exact-match cannot fire.

\subsection{Open-ended generalisation}
\label{sec:method:open}

The same three-component reward (Eq.~\ref{eq:rt}) applies to open-ended
tasks with two substitutions: $v_{\mathrm{match}}$ becomes the continuous
PES judge, and $c$ uses the backward-prediction variant above. No
algorithmic change is required. We show in
Section~\ref{sec:exp:open_ended} that \evoopd{} improves PES scores on
\geneexam{} \texttt{GENE-Arena} by \todo{X} points over vanilla OPD,
demonstrating that the lineage-consistency signal carries the load when
the verifier alone provides only a noisy rubric score.

When $\lambda_v = \lambda_c = 0$, Eq.~\ref{eq:rt} reduces to vanilla
on-policy distillation with per-token role gating; when additionally
$\alpha \equiv 1$, it reduces to the Tinker recipe \citep{tinker_opd_2025}.
This makes \evoopd{} non-inferior to vanilla OPD by construction (modulo
optimisation noise), and the ablations in Sec.~\ref{sec:exp:ablation}
quantify the contribution of each component.

\subsection{Training data shape}
\label{sec:method:data}

Unlike SFT, \evoopd's training data is regenerated each step from
on-policy rollouts. A single training record contains: the prompt $x$
(sampled from a fixed prompt pool of $\sim$10K, drawn from \genetrace{}
+ math/code preservation + \geneexam{} Arena prompts); the student
rollout $y$; the teacher per-token log-probabilities $\log \pi_T(y_t |
y_{<t}, x)$; the per-token role tags $\phi(t)$ from the parser; the
verifier subscores $v(y)$; the lineage anchor and computed $c(y, p)$;
and the resulting per-token reward $r_t$. Full schema in
Appendix~\ref{app:opd_record}.

\section{Experiments}
\label{sec:exp}

\subsection{Setup}
\label{sec:exp:setup}

\paragraph{Student and teacher.} The student is Qwen3-8B
(instruction-tuned) \citep{qwen3_2025}; the OPD teacher is
Qwen3-14B-Thinking (open weights) served locally via vLLM on a single
A100 80GB. All SFT and \evoopd{} training uses LoRA \citep{hu2021lora}
with rank 64, $\alpha=128$, applied to all linear layers.

\paragraph{Benchmarks.}
\geneexam{} \emph{main challenge} (1,029 instances spanning 42 task
types in tiers T1--T4) is the closed-form evaluation. \texttt{GENE-Arena}
(\todo{N} prompts judged on the PES rubric) is the open-ended
evaluation. We additionally report on \texttt{MATH} and \texttt{HumanEval}
to verify no regression on math/code.

\paragraph{Baselines.} (1) Qwen3-8B no-think (baseline); (2) SFT on
\genetrace-derived examples; (3) GRPO with our verifier; (4) Vanilla OPD
(Tinker recipe) with the same teacher; (5) GRPO $\to$ OPD; (6) \evoopd{}.
All variants use matched compute.

\paragraph{Contamination check.} Before training, we compute Min-K\%++
\citep{shi2024minkpp} on \geneexam{} papers against the trained model;
report in Sec.~\ref{sec:exp:contamination} that all variants have
Min-20\%++ $<$ 0.05, consistent with no leakage.

\subsection{\geneexam{} main result}
\label{sec:exp:main}

\todo{Insert Table~\ref{tab:main_geneexam} with strict + lenient
accuracy for all variants. Headline: SFT 11.27\%, \evoopd{} target X\%.}

\begin{table}[t]
  \centering
  \small
  \begin{tabular}{lrr}
    \toprule
    Model            & Strict (\%) & Lenient (\%) \\
    \midrule
    Qwen3-8B (baseline)    & 0.68 & 7.77 \\
    \quad + SFT (\genetrace) & 1.17 & 11.27 \\
    \quad + GRPO            & \todo{} & \todo{} \\
    \quad + vanilla OPD     & \todo{} & \todo{} \\
    \quad + \evoopd{}       & \todo{} & \todo{} \\
    \midrule
    GPT-5.5 (frontier ref.) & 23.10 & 23.10 \\
    \bottomrule
  \end{tabular}
  \caption{\geneexam{} main challenge, no-think mode. Lenient applies
  key-name normalisation; see Appendix~\ref{app:lenient}.}
  \label{tab:main_geneexam}
\end{table}

\subsection{Open-ended (\texttt{GENE-Arena}) result}
\label{sec:exp:open_ended}

\todo{Insert PES scores per dimension; show \evoopd{} > vanilla OPD on
all three (Heredity especially, given the lineage signal).}

\subsection{Component ablations}
\label{sec:exp:ablation}

We ablate the three components of \evoopd{} on \geneexam{} and
\texttt{GENE-Arena}:
\textbf{(L1)} $\alpha \equiv 1$, $\lambda_v=\lambda_c=0$: vanilla OPD;
\textbf{(L2)} drop verifier ($\lambda_v=0$): teacher-KL + lineage only;
\textbf{(L3)} drop lineage ($\lambda_c=0$): teacher-KL + verifier only;
\textbf{(L4)} drop role gating ($\alpha \equiv 1$): all three signals,
no per-token weighting;
\textbf{(L7)} swap Qwen3-14B teacher for GPT-5.5 Black-Box (top-20
logprobs): tests the teacher choice independently.

\todo{Insert Table~\ref{tab:ablation}}.

\subsection{Pipeline ablations}
\label{sec:exp:pipeline}

We ablate the broader training pipeline: SFT-only, SFT $\to$ GRPO,
SFT $\to$ vanilla OPD, SFT $\to$ \evoopd{}, SFT $\to$ GRPO $\to$
\evoopd{}, and \evoopd{} from base (no SFT). This isolates whether
\evoopd{} requires a GRPO warm-up or can be applied directly after SFT.

\subsection{Contamination diagnostic}
\label{sec:exp:contamination}

\todo{Report Min-K\%++ for each \geneexam{} paper subject under each
trained checkpoint; baseline = Qwen3-8B before any fine-tuning.}

\subsection{Cross-domain transfer (v0.2 preview)}
\label{sec:exp:transfer}

To test whether the genome paradigm generalises beyond CS papers, we
evaluate the SFT $\to$ \evoopd{} model on a held-out 200-paper biology
slice from \genetrace{} v0.2. \todo{Report transfer numbers.}

\section{Discussion}
\label{sec:discussion}

\paragraph{What the genome layer adds over method-entity graphs.}
Concurrent work Intern-Atlas demonstrates that 1M-paper method-entity
graphs are constructible automatically. \genetrace{} is much smaller
(5K--50K), but each annotation is structurally richer. Our hypothesis is
that for \emph{training} a model to reason about paper evolution, depth
matters more than breadth --- because the model needs supervision on
\emph{why} edges exist, not just \emph{that} they exist. The
cross-domain transfer experiment (Sec.~\ref{sec:exp:transfer}) is one
test of this hypothesis: if the genome paradigm is genuinely a
representational improvement, models trained on it should transfer
across domains where method-entity vocabulary differs.

\paragraph{When does the lineage-consistency signal help?}
Ablation L2 (drop verifier) and L3 (drop lineage) decompose the
contribution of the two new signals. We expect L3 (no lineage) to lose
the most on open-ended tasks --- where the verifier provides only a
noisy rubric and the lineage signal is the load-bearing reward --- and
L2 (no verifier) to lose the most on closed-form exam tasks. The
asymmetry justifies including both signals rather than picking one.

\paragraph{Teacher staleness during \evoopd{} training.}
On-policy distillation re-queries the teacher every rollout, which is
expensive at scale. We amortise teacher queries by serving the teacher
on a dedicated GPU and batching rollouts; one open question is whether
infrequent teacher updates (every $K$ rollouts) degrade learning, and at
what $K$ the trade-off flips. Sec.~\ref{sec:exp:ablation} L7 probes one
variant of this question.

\paragraph{The contamination guard is a design decision, not just an
audit.} Most paper-graph corpora released to date do not ship with a
denylist or temporal cut explicitly tied to a downstream benchmark. We
argue this should be standard practice: any corpus intended for
training models evaluated on benchmarks derived from the same
literature must demonstrate the disjointness end-to-end, ideally with
re-runnable code. \genetrace{} ships such code with the release.

\paragraph{Failure modes observed during construction.}
During the \genetrace{} v0.1 build, we observed three recurrent failure
modes in teacher annotations: (a) the teacher omits evidence quotes
when the source paper is sparse (e.g., abstract-only); (b) the teacher
prefers \textit{Speciation} over \textit{Mutation} when the paper-pair
$\delta$ is large but the dynamics are actually \textit{Mutation} + new
benchmark; (c) the teacher hallucinates HYBRIDIZED fates for pairs that
share keywords but not actual technique provenance. The verifier
catches (a) and (c) automatically; (b) requires human review on a
sample. We report inter-rater agreement on a 200-instance human-judged
slice in Appendix~\ref{app:human_eval}.


\section*{Limitations}
\genetrace{} v0.1 is restricted to CS papers and to a single teacher
model (GPT-5.5). Cross-domain coverage (v0.2 biology and physics) and
multi-teacher consensus filtering are planned but not in scope for this
release. The six-field genome schema is inherited from our prior work
\ideaevolving{} and was not re-derived from first principles for this
paper; an ablation dropping individual fields (mechanism alone, niche
alone, etc.) would test whether all six are load-bearing. The five-mode
dynamics taxonomy is similarly inherited; alternative taxonomies (e.g.,
finer-grained transition types) are not explored.

\evoopd{} requires a teacher model with accessible per-token
log-probabilities. Open-weights teachers (Qwen3-14B-Thinking,
Llama-3-70B) satisfy this; closed APIs that return only top-$K$
log-probs (GPT-5.5 Black-Box, arXiv 2511.10643) work but with a known
truncation hit. The lineage-consistency signal requires the corpus to
provide multi-paper chains; tasks without natural lineage structure
(single-instance QA, open-domain reasoning) cannot benefit from
component (C) of the algorithm.

The contamination guard relies on the published \ideaevolving{}
\texttt{paper\_db} being a complete enumeration of \geneexam{} source
papers. If \geneexam{} draws from papers outside this list, our
denylist may be incomplete; the temporal cut (pre-2017) provides a
fallback guarantee. We measure Min-K\%++ on the test set as an
independent leakage check.

Reproducibility: all training runs are at LoRA scale and complete in
hours; replicating the full pipeline (SFT + \evoopd{}) requires 4$\times$
A100 80GB for $\sim$3 days. We provide the LoRA adapter weights, the
\genetrace{} v0.1 corpus, and the construction code under permissive
licenses to enable independent replication.


\section*{Ethics Statement}
\paragraph{Source data.} \genetrace{} is annotated over papers from the
\ideaevolving{} \texttt{paper\_db}, which aggregates publicly available
abstracts and intros from ArXiv and conference proceedings. We do not
include full-text PDFs in the release. Paper provenance (\texttt{paper\_id},
year, domain) is preserved per record.

\paragraph{Teacher model outputs.} Genome cards, dynamics labels, and
reasoning traces in \genetrace{} are generated by GPT-5.5 via Azure
OpenAI Service. We are operating under the standard Azure OpenAI
customer terms, which permit derivative use of model outputs for
research and downstream model training. We cite this clause in the
dataset card and pin the API version to enable reproducibility.

\paragraph{Risk of misuse.} A model trained to propose scientific ideas
could be misused to generate plausible but unverified research claims at
scale. We release the model with a card documenting expected use cases
(literature review assistance, idea-generation seed) and explicit
caveats against unsupervised deployment. The verifier and lineage
checks shipped with the release provide a first-line consistency check
on model outputs.

\paragraph{Contamination as an ethics issue.} Training on benchmark
papers and then reporting benchmark scores is a form of evaluation
fraud. Our contamination guard (Sec.~\ref{sec:genetrace:contamination})
is designed to make this category of failure impossible by construction,
and we ship the verification code with the release so reviewers and
downstream consumers can confirm the guarantee independently.

\begin{table}[h]
  \centering
  \small
  \begin{tabular}{lrrr}
    \toprule
    Model & eval$_{\mathrm{mean}}$ & ref$_{\mathrm{mean}}$ & $\Delta$ \\
    \midrule
    Qwen3-8B (untrained)   & +4.1888 & +5.1451 & $-0.9563$ \\
    \quad + SFT v1         & +4.0653 & +4.9478 & $-0.8825$ \\
    \quad + SFT v2         & +4.1049 & +4.9816 & $-0.8767$ \\
    \quad + SFT v3 (best)  & +4.1864 & +5.0593 & $-0.8729$ \\
    \quad + SFT v4         & +4.1950 & +5.0685 & $-0.8735$ \\
    \quad + SFT v5         & +4.1853 & +5.0561 & $-0.8708$ \\
    \quad + SFT v6         & +4.1818 & +5.0525 & $-0.8707$ \\
    \bottomrule
  \end{tabular}
  \caption{Min-20\%++ leakage scores on 50 GENE-bench papers
  (eval) and 50 GeneTrace source papers (ref) per checkpoint.
  $\Delta$ = eval $-$ ref; negative $\Delta$ indicates the trained
  model is \emph{less} confident on eval papers than on
  training-pool papers, the signature of no leakage. All trained
  checkpoints have $\Delta$ within $0.01$ of baseline, well below
  the $0.05$ threshold.}
  \label{tab:leakage}
\end{table}

Table~\ref{tab:leakage} reports Min-20\%++ scores
\citep{shi2024minkpp} for each trained checkpoint. Training on
GeneTrace does not measurably increase model confidence on held-out
\geneexam{} papers (eval$_{\mathrm{mean}}$ for v3 is +4.1864 vs
baseline +4.1888, a difference within the per-paper noise level).

\paragraph{Compute.} Total compute for all reported experiments is
$\sim$\todo{N} A100-hours, dominated by the v0.1 teacher annotation pass
($\sim$\todo{M} A100-hours of GPT-5.5 API time) and the \evoopd{}
training runs ($\sim$\todo{K} A100-hours). We use no specialised
hardware beyond commodity A100 80GB GPUs.


\section*{Acknowledgments}
\todo{Acknowledgments after de-anonymisation.}

\bibliography{custom}

\appendix
\section{GenomeCard JSON Schema}
\label{app:schema:card}

\todo{Insert full JSON schema from \texttt{genetrace\_data\_format.md} §1.}

\section{Verifier Subscores}
\label{app:verifier}

\todo{Insert verifier definitions from \texttt{evo\_opd/verifier.py}.}

\section{evo-OPD Training Record Schema}
\label{app:opd_record}

\todo{Insert OPD record schema from \texttt{evo\_opd\_open\_ended.md} §4.}

\section{Lenient Scoring}
\label{app:lenient}

\todo{Describe key-normalisation dictionary used in lenient evaluation;
list the per-task aliases.}

\section{Human-Eval Inter-Rater Agreement}
\label{app:human_eval}

\todo{200-instance human-judged slice; Cohen's $\kappa$ per dynamics
label.}

\section{Reproducibility Checklist}
\label{app:reproducibility}

\todo{Standard ACL reproducibility checklist.}


\end{document}
